\chapter{Reference: complete example apps}
\label{ch:refman-apps}

\begin{aside}
The previous chapters showed one feature at a time. This
one shows three short but complete programs. Each is small
enough to read in one sitting and exercises several
features at once.
\end{aside}

\section{App 1: ToDo list}

\begin{lstlisting}[language=C++]
#include <maya/maya.hpp>
#include <vector>
#include <string>

using namespace maya;
using namespace maya::dsl;

struct Todo {
    struct Item { std::string text; bool done = false; };
    struct Model {
        std::vector<Item> items;
        std::string buffer;
        int selected = 0;
    };
    struct Type { char c; };
    struct Back {};
    struct Add {};
    struct Toggle {};
    struct Up {};
    struct Down {};
    struct Quit {};
    using Msg = std::variant<Type, Back, Add, Toggle, Up, Down, Quit>;

    static Model init() { return {}; }

    static auto update(Model m, Msg msg)
        -> std::pair<Model, Cmd<Msg>>
    {
        return std::visit(overload{
            [&](Type t) { m.buffer.push_back(t.c); return std::pair{m, Cmd<Msg>{}}; },
            [&](Back) {
                if (!m.buffer.empty()) m.buffer.pop_back();
                return std::pair{m, Cmd<Msg>{}};
            },
            [&](Add) {
                if (!m.buffer.empty()) {
                    m.items.push_back({m.buffer, false});
                    m.buffer.clear();
                }
                return std::pair{m, Cmd<Msg>{}};
            },
            [&](Toggle) {
                if (!m.items.empty() && m.selected < static_cast<int>(m.items.size())) {
                    m.items[m.selected].done = !m.items[m.selected].done;
                }
                return std::pair{m, Cmd<Msg>{}};
            },
            [&](Up) {
                m.selected = std::max(0, m.selected - 1);
                return std::pair{m, Cmd<Msg>{}};
            },
            [&](Down) {
                int n = static_cast<int>(m.items.size());
                m.selected = std::min(std::max(0, n - 1), m.selected + 1);
                return std::pair{m, Cmd<Msg>{}};
            },
            [](Quit) { return std::pair{Model{}, Cmd<Msg>::quit()}; },
        }, msg);
    }

    static Element view(const Model& m) {
        std::vector<Element> rows;
        for (size_t i = 0; i < m.items.size(); ++i) {
            auto& it = m.items[i];
            auto mark = it.done ? "[x] " : "[ ] ";
            auto line = h(text(mark), text(it.text));
            if (static_cast<int>(i) == m.selected) line = line | Bold | Bg<40, 80, 140>;
            if (it.done) line = line | Dim;
            rows.push_back(line);
        }
        return v(
            t<"todo"> | Bold,
            blank_,
            v(rows),
            blank_,
            h(t<"add: ">, text(m.buffer), text("_") | Blink),
            blank_,
            t<"enter add, space toggle, up/down move, q quit"> | Dim
        ) | pad<1> | border_<Round>;
    }

    static auto subscribe(const Model&) -> Sub<Msg> {
        std::vector<std::pair<KeyEvent, Msg>> binds;
        for (char c = ' '; c < 127; ++c) binds.push_back({c, Type{c}});
        binds.push_back({SpecialKey::Backspace, Back{}});
        binds.push_back({SpecialKey::Enter,     Add{}});
        binds.push_back({' ',                   Toggle{}});
        binds.push_back({SpecialKey::Up,        Up{}});
        binds.push_back({SpecialKey::Down,      Down{}});
        binds.push_back({ctrl('c'),             Quit{}});
        return key_map<Msg>(std::move(binds));
    }
};

int main() { run<Todo>({.title = "todo"}); }
\end{lstlisting}

\section{App 2: live system stats}

A version of \code{sysmon} small enough to fit on one page,
suitable as a starting point.

\begin{lstlisting}[language=C++]
#include <maya/maya.hpp>
#include <chrono>
#include <random>

using namespace maya;
using namespace maya::dsl;

struct Stats {
    struct Model {
        std::vector<float> cpu;
        std::vector<float> mem;
        std::mt19937 rng{42};
    };
    struct Tick {};
    struct Quit {};
    using Msg = std::variant<Tick, Quit>;

    static Model init() { return {}; }

    static auto update(Model m, Msg msg)
        -> std::pair<Model, Cmd<Msg>>
    {
        return std::visit(overload{
            [&](Tick) {
                std::uniform_real_distribution<float> d(0.f, 1.f);
                m.cpu.push_back(d(m.rng));
                m.mem.push_back(d(m.rng));
                if (m.cpu.size() > 40) m.cpu.erase(m.cpu.begin());
                if (m.mem.size() > 40) m.mem.erase(m.mem.begin());
                return std::pair{m, Cmd<Msg>{}};
            },
            [](Quit) { return std::pair{Model{}, Cmd<Msg>::quit()}; },
        }, msg);
    }

    static Element spark(const std::vector<float>& data) {
        const char* glyphs[] = {".", ":", "-", "=", "*", "#"};
        std::string s;
        for (float v : data) {
            int idx = std::min(5,
                std::max(0, static_cast<int>(v * 6)));
            s += glyphs[idx];
        }
        return text(s);
    }

    static Element view(const Model& m) {
        return v(
            t<"system">  | Bold,
            blank_,
            h(t<"cpu: ">, spark(m.cpu) | Fg<255, 200, 100>),
            h(t<"mem: ">, spark(m.mem) | Fg<100, 200, 255>),
            blank_,
            t<"q quit"> | Dim
        ) | pad<1> | border_<Round>;
    }

    static auto subscribe(const Model&) -> Sub<Msg> {
        return Sub<Msg>::batch(
            key_map<Msg>({ {'q', Quit{}} }),
            Sub<Msg>::every(std::chrono::milliseconds(200),
                            Tick{})
        );
    }
};

int main() { run<Stats>({.title = "stats"}); }
\end{lstlisting}

\section{Notes}

\begin{itemize}[leftmargin=*]
  \item Both apps fit in under 100 lines.
  \item Both are pure aside from the timer (and the
    todo's buffer-as-state).
  \item Both can be extended without restructuring.
\end{itemize}

\section{Recap}

\begin{recap}
\begin{itemize}[leftmargin=*]
  \item Small apps are a few dozen lines each.
  \item Same skeleton every time.
  \item The grammar is small; the vocabulary is wide.
\end{itemize}
\end{recap}

\section*{Workshop}

\begin{enumerate}[leftmargin=*]
  \item Add a "clear all" key to the todo.
  \item Replace the random source in stats with real \code{/proc} reads.
  \item Combine both into a tab-switched app.
\end{enumerate}
